\begin{exercisechunk}
\item \textbf{Equal margins preserve a signed-gradient ray.}\suppmark
\label[exercise]{ex:l10-equal-margin-ray}
\exerciseprofile{2}{2}{2}{\exproof\quad\excalculation\quad\exinterpretation}{%
Recognize when coordinate normalization confines an adaptive method to one
parameter ray.}
Let $\mathbf X\in\R^{n\times d}$ have row $i$ equal to $x_i^\top$, let
$y\in\{-1,1\}^n$, and define
\[
 s=\operatorname{sign}(\mathbf X^\top y).
\]
Assume $\mathbf Xs=cy$ for some $c>0$.
\begin{enumerate}[(a),beginpenalty=10000]
\item For a differentiable margin loss
$L(w)=\sum_{i=1}^n\ell(y_i\langle w,x_i\rangle)$, prove that
\[
 \nabla L(\alpha s)=\ell'(c\alpha)\mathbf X^\top y.
\]
When $\ell'<0$, prove by induction that signed gradient descent initialized
at zero remains on the ray $\{\alpha s:\alpha\ge0\}$.

\item For the squared loss
$Q(w)=\frac12\|\mathbf Xw-y\|_2^2$, prove
\[
 \nabla Q(\alpha s)=(c\alpha-1)\mathbf X^\top y.
\]
Explain why an adaptive update that divides each nonzero gradient coordinate
by a matching coordinatewise magnitude also remains on the ray.

\item Explain why replacing the denominator $|g_j|$ by $|g_j|+\epsilon$ for
fixed $\epsilon>0$ generally breaks the exact ray calculation as the gradient
vanishes. Contrast the consequences for a finite least-squares interpolant
and for logistic loss on separable data.
\end{enumerate}
\end{exercisechunk}
